\chapter{Conclusion and Perspectives}
%==============================================================================
\pagestyle{fancy}
\fancyhf{}
\fancyhead[R]{\bfseries\rightmark}
\fancyfoot[R]{\thepage}
\renewcommand{\headrulewidth}{0.5pt}
\renewcommand{\footrulewidth}{0pt}
\renewcommand{\chaptermark}[1]{\markboth{\MakeUppercase{\chaptername~\thechapter.\ #1}}{}}
\renewcommand{\sectionmark}[1]{\markright{\thesection~ #1}}

\begin{spacing}{1.2}
%==============================================================================

% TODO: Two parts only (see _Notes/conclusion.md):
%
% PART 1 -- Recap (~1 page)
%   - Contribution: CLMPilot, an integration-first CLM platform for the
%     DACH Mittelstand, built on the comby framework (CQRS + Event
%     Sourcing in Go), shipped in live production at clmpilot.com.
%   - Realised scope: Phases 0-2 complete -- repository foundation, the
%     full-stack MVP (backend domains Contract / Approval / Deadline /
%     DocumentLink, Vue 3 SPA, Hetzner-hosted production deployment with
%     monitoring and backups), and the enterprise tier (intelligence
%     extraction, e-signature, workflow/temporal automation, i18n).
%   - Methodology: release-train delivery with a single source-of-truth
%     roadmap that drove both the work and the structure of this report.
%   - Difficulties encountered (briefly): the .dockerignore /docs gap;
%     the deploy.sh lockstep behaviour; the older comby MFA example
%     bug; auth-vs-query authorisation gap for backfills.
%
% PART 2 -- Perspectives (~1 page)
%   (absorbs the retired Chapter V -- improvements in the event-sourced
%   system -- now grounded in the comby framework audit of 2026-07-07,
%   planning-doc/AUDIT.md; present as an audit performed + evolution
%   plan, at architecture level, no file:line detail)
%   - Evolution of the event-sourced framework: the internship closed
%     with a source-level audit of comby; its findings define the
%     hardening path -- store-enforced optimistic concurrency and atomic
%     event-batch persistence, a corrected snapshot subsystem, a
%     bounded sharded execution model with graceful shutdown, a global
%     sequence with checkpointed projections and a dead-letter queue,
%     and reliable cross-instance replication for read replicas.
%   - Deeper multi-tenancy: per-tenant quotas and metering, physical
%     partitioning of the event store, crypto-shredding for GDPR
%     erasure (from the same audit).
%   - Insight extraction from broker communications (design idea from
%     the retired Chapter V shell): extend the extraction pipeline from
%     contract documents to the communications around a contract.
%   - Phase 3 product roadmap: public API, ERP integrations (DATEV,
%     SAP), multi-entity consolidation, marketplace, batch
%     intelligence, K8s, multi-region.
%   - Beyond the roadmap: open research directions (clause-level
%     reasoning over the event store, federated audit-trail
%     verification, contract similarity embeddings for negotiation
%     support).

%Conclusion

This report presented \textbf{CLMPilot}, an integration-first contract
lifecycle management platform for the DACH \textit{Mittelstand}, built
during an end-of-studies internship at Gradient~Zero and live in
production at \texttt{clmpilot.com}. The platform manages the lifecycle
of a contract (its metadata, deadlines, approvals and audit trail)
while the document itself stays in the customer's own storage, so
adoption carries no migration cost. It is built on the \texttt{comby}
framework, which brings CQRS and Event~Sourcing in Go, and that choice
is what makes the compliance-grade audit trail a property of the
storage model itself: every change to a contract is an immutable event,
traced to an actor and a moment, and any past state can be
reconstructed on demand.

The pilot was delivered as a release train of 3 product phases, all
realised and shipped. Phase~0 laid the engineering foundations:
repositories, continuous integration and delivery, and the roadmap that
served as the single source of truth for the whole internship, ordering
both the work and the structure of this report. Phase~1 delivered the
full-stack MVP: the 4 core domains (Contract, Approval, Deadline and
DocumentLink) behind a generated API, an internationalised web
application, and a monitored, backed-up production deployment on
EU-hosted infrastructure. Phase~2 added the enterprise tier:
intelligence-assisted metadata and clause extraction, electronic
signature, per-tenant custom fields, obligation tracking and advanced
reporting. The work ran under Extreme Programming adapted to a solo
developer, with weekly supervision and release gates in place of team
practices. The main difficulties were operational: delivering alone
against fixed phase milestones, and keeping the releases of 3
separately versioned services reproducible as one unit.

% Perspectives
\vspace{5mm}

The closing weeks of the internship went into an engineering audit of
the event-sourced framework beneath the platform, asking of each
guarantee the architecture depends on where it is enforced and where
that enforcement stops. Most of the perspectives below come from its
findings.

On consistency, the audit found that the framework enforces its
strongest guarantees in the running process and not in the store:
optimistic concurrency is checked against a version cache in memory,
and the events of a command are written one by one. A single instance
behaves correctly, but a second concurrent writer, even a brief
overlap during a deployment, could write conflicting versions into a
stream without any error, and a crash between two writes leaves an
operation half-recorded. The fix is to enforce these rules in the
event store: a uniqueness constraint on aggregate and version, and one
transaction per command. A corrupted stream then becomes a conflict
the dispatcher can catch and retry.

Operating at scale raises the same kind of question. Start-up rebuilds
every read model by replaying the whole log, and the snapshot
mechanism meant to bound that has defects to correct before it can be
enabled. Under load, nothing bounds concurrent work, ordering per
aggregate is not guaranteed, and shutdown can drop accepted commands.
Assigning each aggregate to one worker in a sharded execution model
would give ordering, backpressure and a clean shutdown at once. Read
models also need a recovery story: a global order over the log,
checkpointed projections and a dead-letter queue, so a failing
projection can be repaired and replayed instead of silently skipping
events, and replicas that catch up from the store when broker messages
are missed. This is what it would take to run the multi-instance
topology the deployment design already anticipates.

Isolation between tenants is solid (a tenant context on every command
and event, with row-level security behind it), but growth adds
problems isolation does not solve. Nothing bounds one tenant's load,
so quotas and metering are needed on the shared dispatch path.
Partitioning the event store by tenant would give large or
compliance-sensitive customers their own storage, backup and deletion
boundaries. Corporate groups need consolidation across their tenants,
which has no model yet. And the immutable log conflicts with GDPR
erasure; per-tenant encryption keys would settle that conflict, since
destroying a key erases the personal data while the stream stays
intact.

On the product side, the extraction pipeline is worth extending. It is
a stateless boundary (fetch a document's bytes, extract metadata and
clauses, event-source the result) and nothing in it is specific to
contract documents; the broker communications around a contract, which
the platform does not touch today, could be fed through the same path.
Deeper connections to the storage and ERP systems the
\textit{Mittelstand} already runs would serve the same goal: adoption
that asks no migration of the customer.

The pilot showed that a compliance-grade contract registry can be
delivered at mid-market cost when auditability is built into the
storage model. The work sketched above is what would turn that pilot
into a durable product.

\end{spacing}
